\documentclass{article}

\usepackage{graphicx}
\usepackage{tabularray}
\usepackage{float}
\usepackage{codehigh}
\usepackage[normalem]{ulem}

\UseTblrLibrary{booktabs}
\UseTblrLibrary{siunitx}

\newcommand{\tinytableTabularrayUnderline}[1]{\underline{#1}}
\newcommand{\tinytableTabularrayStrikeout}[1]{\sout{#1}}
\NewTableCommand{\tinytableDefineColor}[3]{\definecolor{#1}{#2}{#3}}

\title{Problem Set 7}
\author{Anjeela Budha Magar}
\date{March 2026}

\begin{document}

\maketitle
\textbf{Part 6. }At what rate are log wages missing? Do you think the logwage variable is most likely to be MCAR, MAR, or MNAR?


Approximately 25\% of observations have missing values for the \texttt{logwage} variable. 

Regarding the type of missingness, the \texttt{logwage} variable is unlikely to be Missing Completely At Random (MCAR), because missing wages could be related to characteristics such as tenure, education, age, and/or employment type. It is more plausible that the missingness is Missing At Random (MAR), meaning that the probability of \texttt{logwage} being missing depends on observed variables in the dataset such as education or age. Missing Not At Random (MNAR) is less likely in this case, but cannot be fully ruled out without additional information on why wages are missing.

\begin{table}
\centering
\begin{tblr}[         %% tabularray outer open
]                     %% tabularray outer close
{                     %% tabularray inner open
colspec={Q[]Q[]Q[]Q[]Q[]Q[]Q[]Q[]Q[]},
hline{2}={1-9}{solid, black, 0.05em},
hline{1}={1-9}{solid, black, 0.08em},
hline{11}={1-9}{solid, black, 0.08em},
hline{6}={1-9}{solid, black, 0.1em},
}                     %% tabularray inner close
& Unique & Missing Pct. & Mean & SD & Min & Median & Max & Histogram \\
logwage & 670 & 25 & 1.6 & 0.4 & 0.0 & 1.7 & 2.3 &  \\
hgc & 16 & 0 & 13.1 & 2.5 & 0.0 & 12.0 & 18.0 &  \\
tenure & 259 & 0 & 6.0 & 5.5 & 0.0 & 3.8 & 25.9 &  \\
age & 13 & 0 & 39.2 & 3.1 & 34.0 & 39.0 & 46.0 &  \\
&  & N & % &   &   &   &   &   \\
college & college grad & 530 & 23.8 &   &   &   &   &   \\
& not college grad & 1699 & 76.2 &   &   &   &   &   \\
married & married & 1431 & 64.2 &   &   &   &   &   \\
& single & 798 & 35.8 &   &   &   &   &   \\
\end{tblr}
\end{table}
\newpage

\textbf{Part 7. }The true value of  $\hat{\beta}_1$ is 0.093. Comment on the differences of $\hat{\beta}_1$ across the models. What patterns do you see? What can you conclude about the veracity of the various imputation methods? Also discuss what the estimates of $\hat{\beta}_1$ are for the last two methods.

The true value of $\hat{\beta}_1$ is 0.093.  The complete case estimate of $\hat{\beta}_1$ is approximately 0.062, while the mean imputation estimate is lower at around 0.050. The regression imputation method produces an estimate similar to the complete case result, around 0.062. The multiple imputation estimate is approximately 0.060.

All four methods produce estimates of $\hat{\beta}_1$ that are lower than the true value of 0.093, indicating a downward bias across. Mean imputation yields the smallest estimate, which is expected because it reduces the variability in the data and biases coefficients toward zero. Regression imputation produces more stable estimates but may overstate precision since it does not fully account for uncertainty in the imputed values.

The multiple imputation method provides a more reliable estimate compared to the other approaches because it incorporates uncertainty from the imputation process by averaging across multiple datasets. Although its estimate (0.060) is still below the true value, it is more credible than the alternatives in terms of statistical validity.

Multiple imputation is typically preferred, even though some bias may still remain. The consistency between the complete case and the regression imputation estimates also suggests that missingness may not be completely random, reinforcing the earlier conclusion that the data are likely Missing At Random (MAR).
\section*{Regression Results}

\begin{table}[H]
\centering
\begin{tblr}{
colspec={Q[]Q[]Q[]Q[]Q[]},
hline{1}={1-5}{solid, black, 0.08em},
hline{2}={1-5}{solid, black, 0.05em},
hline{16}={1-5}{solid, black, 0.05em},
hline{24}={1-5}{solid, black, 0.08em},
column{2-5}={c},
column{1}={l},
}
& Complete Case & Mean Imputation & Regression Imputation & Multiple Imputation \\
(Intercept) & \num{0.534} & \num{0.708} & \num{0.534} & \num{0.593} \\
& (\num{0.146}) & (\num{0.116}) & (\num{0.112}) & (\num{0.138}) \\
hgc & \num{0.062} & \num{0.050} & \num{0.062} & \num{0.060} \\
& (\num{0.005}) & (\num{0.004}) & (\num{0.004}) & (\num{0.005}) \\
collegenot college grad & \num{0.145} & \num{0.168} & \num{0.145} & \num{0.121} \\
& (\num{0.034}) & (\num{0.026}) & (\num{0.025}) & (\num{0.035}) \\
tenure & \num{0.050} & \num{0.038} & \num{0.050} & \num{0.044} \\
& (\num{0.005}) & (\num{0.004}) & (\num{0.004}) & (\num{0.005}) \\
I(tenure\textasciicircum{}2) & \num{-0.002} & \num{-0.001} & \num{-0.002} & \num{-0.001} \\
& (\num{0.000}) & (\num{0.000}) & (\num{0.000}) & (\num{0.000}) \\
age & \num{0.000} & \num{0.000} & \num{0.000} & \num{0.000} \\
& (\num{0.003}) & (\num{0.002}) & (\num{0.002}) & (\num{0.003}) \\
marriedsingle & \num{-0.022} & \num{-0.027} & \num{-0.022} & \num{-0.015} \\
& (\num{0.018}) & (\num{0.014}) & (\num{0.013}) & (\num{0.017}) \\
Num.Obs. & \num{1669} & \num{2229} & \num{2229} & \num{2229} \\
Num.Imp. &  &  &  & \num{5} \\
R2 & \num{0.208} & \num{0.147} & \num{0.277} & \num{0.228} \\
R2 Adj. & \num{0.206} & \num{0.145} & \num{0.275} & \num{0.226} \\
AIC & \num{1179.9} & \num{1091.2} & \num{925.5} &  \\
BIC & \num{1223.2} & \num{1136.8} & \num{971.1} &  \\
Log.Lik. & \num{-581.936} & \num{-537.580} & \num{-454.737} &  \\
RMSE & \num{0.34} & \num{0.31} & \num{0.30} &  \\
\end{tblr}
\caption{Regression results under different imputation methods}
\end{table}

\textbf{Part 8. }Tell me about the progress you’ve made on your project. What data are you using? What kinds of modeling approaches do you think you’re going to take?

The datasets I might include are GDP per Capita, GDP growth rate, proportion of seats held by women in national parliaments of the World Bank, Women Mobility Score, and Female Labor Force Participation. They are from World Bank Group.\\
So far, I have planned two models:
\vspace{0.5em}

\textbf{GDP growth as a function of FLFP and time:}  
\[
\text{GDP Growth}_{it} = \beta_0 + \beta_1 \text{FLFP}_{it} + \beta_2 \text{Year}_t + \epsilon_{it}
\]
\vspace{0.5em}
\textbf{Interaction with time to assess changing effects:}  
\[
\text{GDP Growth}_{it} = \beta_0 + \beta_1 \text{FLFP}_{it} + \beta_2 \text{Year}_t + \beta_3 (\text{FLFP}_{it} \times \text{Year}_t) + \epsilon_{it}
\]

\end{document}